\section{Loss Function Design}
\label{sec:loss}

\subsection{Displacement Regression Losses: Why They Fail}

The naive approach is to train the network to minimise the L1 (or L2) distance
between predicted displacements and precomputed nearest-neighbour targets:
\begin{equation}
  \mathcal{L}_{\text{L1}} = \frac{1}{N}\sum_{i=1}^{N}
    \left\|\Delta_i - d_{gt}(\hat{p}_i)\right\|_1,
  \qquad d_{gt}(\hat{p}_i) = p^*_{nn}(\hat{p}_i) - \hat{p}_i.
\end{equation}

This fails due to two structural problems:

\paragraph{Problem 1: Zero-inflation.}
At r2, 60\% of decoded points have $d_{gt}(\hat{p}_i) = \mathbf{0}$
(already at a correct voxel position).
The optimal constant predictor that minimises the L1 loss is the coordinate-wise
median of the targets.
When the median of 60\% of targets is zero, the optimal constant prediction
is $\mathbf{0}$ for each coordinate.
The gradient signal from the 40\% non-zero targets is overwhelmed by the 60\%
zero targets pulling the prediction toward zero.
In practice, any network trained with this loss collapses to predicting near-zero
for all points --- the \emph{zero-inflation failure mode}.

We verified this empirically: networks trained with L1, smooth-L1, and
curvature-weighted L1 all collapse to near-zero predictions within 5--10 epochs,
producing negative or negligible PSNR improvement (Table~\ref{tab:ablation-loss}).
Figure~\ref{fig:zero-inflation} illustrates the zero-inflation problem and how
it worsens with rate: at r6, 87\% of decoded points have zero displacement
target, leaving almost no learnable signal.

\begin{figure}[t]
  \centering
  \includegraphics[width=0.88\textwidth]{zero_inflation.pdf}
  \caption{%
    Displacement signal statistics at training rates.
    \textbf{Left}: fraction of decoded points with ground-truth displacement
    $\|d_{gt}\| = 0$ (already correctly placed).
    \textbf{Right}: mean displacement magnitude.
    At r2 (the richest training rate), 60\% of targets are zero and the
    mean displacement is only 0.15 voxels.
    At r6, 87\% zero inflation leaves almost no gradient signal for
    displacement losses.
  }
  \label{fig:zero-inflation}
\end{figure}

\paragraph{Problem 2: Nearest-neighbour noise.}
Even the 40\% non-zero targets are unreliable at edge regions.
At a decoded point $\hat{p}$ near a sharp feature, the nearest neighbour in
$\pc^*$ is not well-defined: a small perturbation of $\hat{p}$ can flip its
nearest neighbour from one side of the edge to the other, producing a
discontinuous jump in $d_{gt}$.
This instability is most severe at exactly the locations where correction
is most needed (high-curvature regions), causing the gradient signal to be
noisiest where it should be most informative.

\subsection{Dynamic Chamfer Distance}
\label{ssec:dcd}

We resolve both problems by abandoning precomputed per-point targets entirely
and directly minimising the geometric discrepancy between the refined patch
and the corresponding crop of the original cloud.

\paragraph{Patch extraction.}
The point cloud is divided into overlapping cubic patches of $64^3$ voxels
with stride 32.
For each patch of decoded points $\hat{\pc}_{\text{patch}} \subset \hat{\pc}$,
the corresponding crop of the original cloud is:
\begin{equation}
  \pc^*_{\text{crop}} = \{p^* \in \pc^* : p^* \in [u_{\min},\, u_{\max}]^3\},
\end{equation}
where $[u_{\min}, u_{\max}]^3$ is the axis-aligned bounding box of
$\hat{\pc}_{\text{patch}}$ with \emph{zero} padding.

\paragraph{Loss computation.}
Let $\hat{\pc}_{\text{ref}} = \{\hat{p}_i + \Delta_i : \hat{p}_i \in \hat{\pc}_{\text{patch}}\}$
be the refined patch.
The per-patch loss is the symmetric Chamfer Distance:
\begin{equation}
  \mathcal{L}_{\text{CD}} =
    \frac{1}{|\hat{\pc}_{\text{ref}}|}\sum_{\hat{p} \in \hat{\pc}_{\text{ref}}}
      \min_{p^* \in \pc^*_{\text{crop}}} \|\hat{p} - p^*\|^2
    + \frac{1}{|\pc^*_{\text{crop}}|}\sum_{p^* \in \pc^*_{\text{crop}}}
      \min_{\hat{p} \in \hat{\pc}_{\text{ref}}} \|p^* - \hat{p}\|^2.
  \label{eq:cd-loss}
\end{equation}

Figure~\ref{fig:chamfer-diagram} illustrates the two terms of the symmetric
Chamfer Distance and why symmetry is essential.

\begin{figure}[t]
\centering
\begin{tikzpicture}[>=Stealth, every node/.style={font=\small}, scale=0.95]
  % ---- Forward term: refined -> original ----
  \begin{scope}
    \node[font=\small\bfseries] at (2.2, 3.1) {Forward: $\hat{\mathcal{P}}_\text{ref} \to \mathcal{P}^*$};
    % Original cloud (green dots, curve shape)
    \foreach \x/\y in {0.5/1.5, 1.0/2.0, 1.5/2.4, 2.0/2.5, 2.5/2.4,
                        3.0/2.0, 3.5/1.5, 4.0/1.2} {
      \fill[green!60!black] (\x,\y) circle (3pt);
    }
    % Refined points (blue dots, slightly off)
    \foreach \x/\y in {0.6/1.0, 1.2/1.5, 1.8/1.9, 2.4/1.8, 3.1/1.4} {
      \fill[blue!70] (\x,\y) circle (3pt);
    }
    % Forward arrows (blue to nearest green)
    \draw[->, red!70, dashed] (0.6,1.0) -- (0.5,1.5);
    \draw[->, red!70, dashed] (1.2,1.5) -- (1.0,2.0);
    \draw[->, red!70, dashed] (1.8,1.9) -- (1.5,2.4);
    \draw[->, red!70, dashed] (2.4,1.8) -- (2.5,2.4);
    \draw[->, red!70, dashed] (3.1,1.4) -- (3.0,2.0);
  \end{scope}

  % ---- Reverse term: original -> refined ----
  \begin{scope}[xshift=6.5cm]
    \node[font=\small\bfseries] at (2.2, 3.1) {Reverse: $\mathcal{P}^* \to \hat{\mathcal{P}}_\text{ref}$};
    % Original cloud
    \foreach \x/\y in {0.5/1.5, 1.0/2.0, 1.5/2.4, 2.0/2.5, 2.5/2.4,
                        3.0/2.0, 3.5/1.5, 4.0/1.2} {
      \fill[green!60!black] (\x,\y) circle (3pt);
    }
    % Refined points
    \foreach \x/\y in {0.6/1.0, 1.2/1.5, 1.8/1.9, 2.4/1.8, 3.1/1.4} {
      \fill[blue!70] (\x,\y) circle (3pt);
    }
    % Reverse arrows (green to nearest blue) — note: some green points far from any blue
    \draw[->, orange!80, dashed] (0.5,1.5) -- (0.6,1.0);
    \draw[->, orange!80, dashed] (1.0,2.0) -- (1.2,1.5);
    \draw[->, orange!80, dashed] (1.5,2.4) -- (1.8,1.9);
    \draw[->, orange!80, dashed] (2.0,2.5) -- (1.8,1.9);
    \draw[->, orange!80, dashed] (2.5,2.4) -- (2.4,1.8);
    \draw[->, orange!80, dashed] (3.0,2.0) -- (3.1,1.4);
    \draw[->, orange!80, thick, dashed] (3.5,1.5) -- (3.1,1.4);
    \draw[->, orange!80, thick, dashed] (4.0,1.2) -- (3.1,1.4);
    % Uncovered region annotation
    \draw[red, thick] (3.3,0.9) ellipse (0.9cm and 0.35cm);
    \node[red, font=\scriptsize, right] at (4.2, 0.9) {uncovered};
  \end{scope}

  % Shared legend -- centered below both panels (global coords)
  \fill[green!60!black] (1.5, -0.35) circle (3pt);
  \node[right, font=\small] at (1.65, -0.35) {Original $\mathcal{P}^*$};
  \fill[blue!70] (4.8, -0.35) circle (3pt);
  \node[right, font=\small] at (4.95, -0.35) {Refined $\hat{\mathcal{P}}_\text{ref}$};
  \draw[->, red!70, dashed] (8.2,-0.45) -- (8.7,-0.25);
  \node[right, font=\small] at (8.75,-0.35) {NN match};
\end{tikzpicture}
\caption{%
  Symmetric Chamfer Distance.
  \textbf{Left (forward term)}: each refined point finds its nearest original point.
  Penalises refined points that are far from the surface.
  \textbf{Right (reverse term)}: each original point finds its nearest refined point.
  Penalises surface regions not covered by the refined cloud (circled).
  Together, both terms enforce that the refined cloud is close to and covers
  the original surface.
}
\label{fig:chamfer-diagram}
\end{figure}

\paragraph{Why this resolves zero-inflation.}
The loss in Equation~\ref{eq:cd-loss} is not computed per-point against a
fixed precomputed target; it is computed over the entire refined set.
For a point $\hat{p}_i$ that is already correctly placed (optimal
$\Delta_i = \mathbf{0}$), its nearest-neighbour distance to $\pc^*_{\text{crop}}$
is already minimal, and the gradient contribution from this point is near zero.
The gradient is therefore concentrated on the \emph{incorrectly placed}
points --- those that are far from their nearest counterpart in $\pc^*_{\text{crop}}$ ---
regardless of what fraction of points are zero.
Zero-inflation does not bias the gradient.

\paragraph{Why this resolves nearest-neighbour noise.}
In Equation~\ref{eq:cd-loss}, the nearest-neighbour matching is performed
on the \emph{refined} positions $\hat{p}_i + \Delta_i$, not the fixed
decoded positions.
As the network updates $\Delta_i$ during training, the matching adapts
accordingly.
This is the ``dynamic'' aspect of the loss: matching is re-computed at
every training step.
The instability of 1-NN at edge regions does not propagate because the
network is free to move points to any position within $[-5, 5]^3$ and the
matching will find the globally optimal nearest neighbour at each step.

\paragraph{Critical implementation detail: zero boundary padding.}
The crop $\pc^*_{\text{crop}}$ must be extracted with \emph{exactly the same}
bounding box as $\hat{\pc}_{\text{patch}}$, i.e., zero padding.

To see why, consider what happens with non-zero padding of $d$ voxels:
the bounding box of $\pc^*_{\text{crop}}$ is expanded by $d$ on each side.
Boundary decoded points near the edge of the patch now have original points
$d$ voxels \emph{outside} the patch, making the reverse term
$\pc^*_{\text{crop}} \to \hat{\pc}_{\text{ref}}$ small (close to 0) regardless
of predictions.
The reverse term becomes negligible and the forward term
$\hat{\pc}_{\text{ref}} \to \pc^*_{\text{crop}}$ dominates.
This creates a gradient that pushes decoded points \emph{outward toward the
patch boundary} (to get closer to the boundary-adjacent original points),
a degenerate signal.

We verified this empirically: with $d = 10$, the reverse term constitutes
over 99\% of the total loss by epoch 3, and the model displaces all points
toward the patch boundaries regardless of rate.
With $d = 0$, the two terms are balanced and the model learns correct
geometric corrections.

% ============================================================
